% GENERATED FILE. Edit canonical lessons, then run npm run generate:books.
% canonical-source-hash: fnv1a64:dd49bfa7de3684a7
% canonical-lessons: SA-C441-durdinam, SA-C441-sajalah, SA-C441-ativrstih, SA-C441-tuhinam, SA-C441-avrnoti

\chapter{Bad Weather}
\label{ch:sa-bad-weather}
\hlchaptermodality{441}

\begin{chapteropening}
\textbf{By the end of this chapter:} \emph{I can say whether the weather is fit for travel.}

Complete the last lesson of chapter 441: I can say whether the weather is fit for travel.
\end{chapteropening}

\section[\texorpdfstring{\sk{दुर्दिनम्} (durdinam)}{दुर्दिनम् (durdinam)}]{\sk{दुर्दिनम्} (durdinam) — a day of bad weather}

\label{lesson:SA-C441-durdinam}



\begin{quote}
Before the new one: say the Sanskrit for huddles up, then the Sanskrit for a cloak.
\end{quote}

\subsection*{You'll want to know: \sk{दुर्दिनम्}}
\textbf{\sk{दुर्दिनम्}} — \emph{durdinam} — \textquotedblleft{}a day of bad weather\textquotedblright{}.

\begin{quote}
\sk{दुर्दिने} \sk{यात्रा} \sk{कठिना} \sk{भवति।}
\end{quote}

\begin{quote}
\emph{durdine yātrā kaṭhinā bhavati.}
\end{quote}

\begin{quote}
On a day of bad weather a journey is hard.
\end{quote}

A neuter noun.

\begin{cousinweb}
\textbf{\sk{दुर्दिनम्}} (\emph{durdinam}) is built from \emph{dur}, \textquotedblleft{}bad, ill\textquotedblright{}, and \emph{dina}, \textquotedblleft{}a day\textquotedblright{}; MW gives the compound as a rainy or cloudy day, bad weather (MW).
\end{cousinweb}

\begin{grammarlens}[title={What you've built}]
One more word for getting around.
\end{grammarlens}

\subsection*{Your turn}
\begin{itemize}
\raggedright
  \item \emph{Say it:} \emph{durdinam}
  \item \emph{Say it:} \emph{durdinam}, once more
  \item \emph{Say it:} \emph{durdine yātrā kaṭhinā bhavati.} — On a day of bad weather a journey is hard.
  \item \emph{Recall:} say the Sanskrit for huddles up, then the Sanskrit for a cloak, then say \emph{durdinam} again
\end{itemize}

\begin{tcolorbox}[breakable,colback=teal!4,colframe=teal!35!black,title={Before you move on}]
What does \sk{दुर्दिनम्} mean? (\textquotedblleft{}A day of bad weather\textquotedblright{}.) Say it once more.
\end{tcolorbox}
\section[\texorpdfstring{\sk{सजलः} (sajalaḥ)}{सजलः (sajalaḥ)}]{\sk{सजलः} (sajalaḥ) — rain-laden; watery}

\label{lesson:SA-C441-sajalah}



\begin{quote}
Before the new one: say the Sanskrit for a cloak, then the Sanskrit for a day of bad weather.
\end{quote}

\subsection*{You'll want to know: \sk{सजलः}}
\textbf{\sk{सजलः}} — \emph{sajalaḥ} — \textquotedblleft{}rain-laden; watery\textquotedblright{}.

\begin{quote}
\sk{सजलम्} \sk{अभ्रं} \sk{पर्वतं} \sk{प्रति} \sk{गच्छति।}
\end{quote}

\begin{quote}
\emph{sajalam abhraṃ parvataṃ prati gacchati.}
\end{quote}

\begin{quote}
A rain-laden cloud moves towards the mountain.
\end{quote}

\begin{cousinweb}
\textbf{\sk{सजलः}} (\emph{sajalaḥ}) is built from \emph{sa}, \textquotedblleft{}with\textquotedblright{}, and \emph{jala}, \textquotedblleft{}water\textquotedblright{}; MW gives \emph{sajala} as containing water, watery, wet (MW).
\end{cousinweb}

\begin{grammarlens}[title={What you've built}]
One more word for getting around.
\end{grammarlens}

\subsection*{Your turn}
\begin{itemize}
\raggedright
  \item \emph{Say it:} \emph{sajalaḥ}
  \item \emph{Say it:} \emph{sajalaḥ}, once more
  \item \emph{Say it:} \emph{sajalam abhraṃ parvataṃ prati gacchati.} — A rain-laden cloud moves towards the mountain.
  \item \emph{Recall:} say the Sanskrit for a cloak, then the Sanskrit for a day of bad weather, then say \emph{sajalaḥ} again
\end{itemize}

\begin{tcolorbox}[breakable,colback=teal!4,colframe=teal!35!black,title={Before you move on}]
What does \sk{सजलः} mean? (\textquotedblleft{}Rain-laden; watery\textquotedblright{}.) Say it once more.
\end{tcolorbox}
\section[\texorpdfstring{\sk{अतिवृष्टिः} (ativṛṣṭiḥ)}{अतिवृष्टिः (ativṛṣṭiḥ)}]{\sk{अतिवृष्टिः} (ativṛṣṭiḥ) — too much rain}

\label{lesson:SA-C441-ativrstih}



\begin{quote}
Before the new one: say the Sanskrit for a day of bad weather, then the Sanskrit for rain-laden.
\end{quote}

\subsection*{You'll want to know: \sk{अतिवृष्टिः}}
\textbf{\sk{अतिवृष्टिः}} — \emph{ativṛṣṭiḥ} — \textquotedblleft{}too much rain\textquotedblright{}.

\begin{quote}
\sk{अतिवृष्ट्या} \sk{क्षेत्राणि} \sk{नश्यन्ति।}
\end{quote}

\begin{quote}
\emph{ativṛṣṭyā kṣetrāṇi naśyanti.}
\end{quote}

\begin{quote}
Too much rain ruins the fields.
\end{quote}

A feminine noun.

\begin{cousinweb}
\textbf{\sk{अतिवृष्टिः}} (\emph{ativṛṣṭiḥ}) is built from \emph{ati}, \textquotedblleft{}beyond, too much\textquotedblright{}, and \emph{vṛṣṭi}, \textquotedblleft{}rain\textquotedblright{}, on the root √vṛṣ; MW gives the compound as excessive rain (MW).
\end{cousinweb}

\begin{grammarlens}[title={What you've built}]
One more word for getting around.
\end{grammarlens}

\subsection*{Your turn}
\begin{itemize}
\raggedright
  \item \emph{Say it:} \emph{ativṛṣṭiḥ}
  \item \emph{Say it:} \emph{ativṛṣṭiḥ}, once more
  \item \emph{Say it:} \emph{ativṛṣṭyā kṣetrāṇi naśyanti.} — Too much rain ruins the fields.
  \item \emph{Recall:} say the Sanskrit for a day of bad weather, then the Sanskrit for rain-laden, then say \emph{ativṛṣṭiḥ} again
\end{itemize}

\begin{tcolorbox}[breakable,colback=teal!4,colframe=teal!35!black,title={Before you move on}]
What does \sk{अतिवृष्टिः} mean? (\textquotedblleft{}Too much rain\textquotedblright{}.) Say it once more.
\end{tcolorbox}
\section[\texorpdfstring{\sk{तुहिनम्} (tuhinam)}{तुहिनम् (tuhinam)}]{\sk{तुहिनम्} (tuhinam) — frost; hoarfrost}

\label{lesson:SA-C441-tuhinam}



\begin{quote}
Before the new one: say the Sanskrit for rain-laden, then the Sanskrit for too much rain.
\end{quote}

\subsection*{You'll want to know: \sk{तुहिनम्}}
\textbf{\sk{तुहिनम्}} — \emph{tuhinam} — \textquotedblleft{}frost; hoarfrost\textquotedblright{}.

\begin{quote}
\sk{शिशिरे} \sk{प्रातः} \sk{तृणेषु} \sk{तुहिनम्} \sk{अस्ति।}
\end{quote}

\begin{quote}
\emph{śiśire prātaḥ tṛṇeṣu tuhinam asti.}
\end{quote}

\begin{quote}
In late winter there is frost on the grass in the morning.
\end{quote}

A neuter noun.

\begin{cousinweb}
\textbf{\sk{तुहिनम्}} (\emph{tuhinam}) is the word MW gives for frost, cold, mist, dew and snow; MW gives no derivation, so its origin is uncertain (MW).
\end{cousinweb}

\begin{grammarlens}[title={What you've built}]
One more word for getting around.
\end{grammarlens}

\subsection*{Your turn}
\begin{itemize}
\raggedright
  \item \emph{Say it:} \emph{tuhinam}
  \item \emph{Say it:} \emph{tuhinam}, once more
  \item \emph{Say it:} \emph{śiśire prātaḥ tṛṇeṣu tuhinam asti.} — In late winter there is frost on the grass in the morning.
  \item \emph{Recall:} say the Sanskrit for rain-laden, then the Sanskrit for too much rain, then say \emph{tuhinam} again
\end{itemize}

\begin{tcolorbox}[breakable,colback=teal!4,colframe=teal!35!black,title={Before you move on}]
What does \sk{तुहिनम्} mean? (\textquotedblleft{}Frost; hoarfrost\textquotedblright{}.) Say it once more.
\end{tcolorbox}
\section[\texorpdfstring{\sk{आवृणोति} (āvṛṇoti)}{आवृणोति (āvṛṇoti)}]{\sk{आवृणोति} (āvṛṇoti) — veils; hides from view}

\label{lesson:SA-C441-avrnoti}



\begin{quote}
Before the new one: say the Sanskrit for too much rain, then the Sanskrit for hoarfrost.
\end{quote}

\subsection*{You'll want to know: \sk{आवृणोति}}
\textbf{\sk{आवृणोति}} — \emph{āvṛṇoti} — \textquotedblleft{}veils; hides from view\textquotedblright{}.

\begin{quote}
\sk{अभ्रं} \sk{सूर्यम्} \sk{आवृणोति।}
\end{quote}

\begin{quote}
\emph{abhraṃ sūryam āvṛṇoti.}
\end{quote}

\begin{quote}
A cloud veils the sun.
\end{quote}

Past (it …): \sk{अभ्रं} \sk{सूर्यम्} \sk{आवृतवत्}, \emph{abhraṃ sūryam āvṛtavat}.

\begin{cousinweb}
\textbf{\sk{आवृणोति}} (\emph{āvṛṇoti}) is built from \emph{ā} and the root √vṛ, \textquotedblleft{}to cover\textquotedblright{}; MW gives \emph{ā-vṛ} as to cover, hide, conceal, with the participle \emph{āvṛta}, \textquotedblleft{}covered\textquotedblright{} (MW).
\end{cousinweb}

\begin{grammarlens}[title={What you've built}]
One more word for getting around.
\end{grammarlens}

\subsection*{Your turn}
\begin{itemize}
\raggedright
  \item \emph{Say it:} \emph{āvṛṇoti}
  \item \emph{Say it:} \emph{āvṛṇoti}, once more
  \item \emph{Say it:} \emph{abhraṃ sūryam āvṛṇoti.} — A cloud veils the sun.
  \item \emph{Recall:} say the Sanskrit for too much rain, then the Sanskrit for hoarfrost, then say \emph{āvṛṇoti} again
\end{itemize}

\begin{tcolorbox}[breakable,colback=teal!4,colframe=teal!35!black,title={Before you move on}]
What does \sk{आवृणोति} mean? (\textquotedblleft{}Veils; hides from view\textquotedblright{}.) Say it once more.
\end{tcolorbox}
